\scanpage{116}
Soit $(r_\nu)$ suite croissante de nombres \deleted{\unsure{positifs}} $>0$, $\to\infty$
soit $n(r)$ le nombre des indices $\nu$ tels que
$r_\nu\leq r$, on a, pour $0\leq a\leq b<+\infty$, et $\alpha>0$
\[ \hbox{\begingroup\setbox0=\hbox{$\displaystyle \sum_{\substack{\nu\\(a<r_\nu<b)}}\frac1{r_\nu^\alpha}=\left(\frac{n(b-0)}{b^\alpha}-\frac{n(a+0)}{a^\alpha}\right)+\alpha\int_a^b\frac{n(r)\,dr}{r^{\alpha+1}}$}\kern5pt\rlap{\tikz[baseline=0pt]{\draw[line width=.35pt] (-1pt,\the\dimexpr\ht0+4pt\relax)--(\the\dimexpr\wd0+7pt\relax,\the\dimexpr\ht0+4pt\relax);\draw[line width=.35pt] (-5pt,\the\dimexpr\ht0+1pt\relax)--(-5pt,\the\dimexpr-\dp0-1pt\relax);\draw[line width=.45pt] (\the\dimexpr\wd0+5pt\relax,\the\dimexpr\ht0+1pt\relax)--(\the\dimexpr\wd0+5pt\relax,\the\dimexpr-\dp0-1pt\relax);\draw[line width=.4pt] (-1pt,\the\dimexpr-\dp0-4pt\relax)--(\the\dimexpr\wd0+5pt\relax,\the\dimexpr-\dp0-4pt\relax);}}\copy0\kern7pt\endgroup}\tag*{(1)} \]
(car le premier membre est égal à : $\int_{\unsure{r_0|a_0}}^b r^{-\alpha}\,dn(r)$,
et on intègre par parties)

\uline{Proposition} \quad $\displaystyle \sum_\nu\frac1{r_\nu^\alpha}<+\infty\ \Longleftrightarrow\ \int_0^\infty\frac{n(r)\,dr}{r^{\alpha+1}}\leq+\infty$
\margin{\tikz[x=1pt,y=1pt]{\draw[line width=.35pt] (0,0)--(.4,-68);}}
alors $\frac{n(r)}{r^\alpha}\to0$ pour $r\to\infty$, et on a
\[ \hbox{\begingroup\setbox0=\hbox{$\displaystyle \sum_{\substack{\nu\\(r_\nu>a)}}\frac1{r_\nu}=-\frac{n(a+0)}{a^\alpha}+\alpha\int_a^\infty\frac{n(r)\,dr}{r^{\alpha+1}}$}\kern5pt\rlap{\tikz[baseline=0pt]{\draw[line width=.35pt] (-1pt,\the\dimexpr\ht0+4pt\relax)--(\the\dimexpr\wd0+7pt\relax,\the\dimexpr\ht0+4pt\relax);\draw[line width=.35pt] (-5pt,\the\dimexpr\ht0+1pt\relax)--(-5pt,\the\dimexpr-\dp0-1pt\relax);\draw[line width=.45pt] (\the\dimexpr\wd0+5pt\relax,\the\dimexpr\ht0+1pt\relax)--(\the\dimexpr\wd0+5pt\relax,\the\dimexpr-\dp0-1pt\relax);\draw[line width=.4pt] (2pt,\the\dimexpr-\dp0-4pt\relax)--(\the\dimexpr\wd0+5pt\relax,\the\dimexpr-\dp0-4pt\relax);}}\copy0\kern7pt\endgroup}\tag*{(1)} \]

$\leftarrow$
\deleted{\unsure{si} $\sum_\nu\frac1{r_\nu^\alpha}<+\infty$, \unsure{on a} $\frac{n(r)}{r^\alpha}\leq M$ pour \deleted{\unsure{tout}} $r$
où $M=$ \unsure{constante}}
En effet, (1) s'écrit, pour $a=0$, \deleted{\unsure{comme|donc}}
\[ \sum_{r_\nu<b}\frac1{r_\nu^{\unsure{\nu|\alpha}}}=\frac{n(b-0)}{b^\alpha}+\alpha\int_0^b\frac{n(r)\,dr}{r^{\alpha+1}}\tag*{\textcircled{\scriptsize 1 bis}} \]
d'où suit que si \deleted{\unsure{pour un|pour} $b$} $\sum_\nu\frac1{r_\nu^\alpha}$ est
convergente, il en est de $=$ de $\int_0^\infty\frac{n(r)\,dr}{r^{\alpha+1}}$
Réciproquement, si $\int_{\unsure{0}}^\infty\frac{n(r)\,dr}{r^{\alpha+1}}<+\infty$ alors il
existe $M>0$ et une suite croiss. de $b_i\to\infty$
de nbs tels que $\frac{n(b_i)}{b_i^\alpha}\leq M$ ; \unsure{p.} on a donc
$\displaystyle \sum_{r_\nu<b_i}\frac1{r_\nu^\alpha}\leq M+\alpha\int_0^\infty\frac{n(r)\,dr}{r^{\alpha+1}}$, d'où $\displaystyle \hbox{\begingroup\setbox0=\hbox{$\displaystyle\sum_{\unsure{\nu|i}}\frac1{r_{\unsure{\nu|i}}^\alpha}$}\rlap{\raisebox{\dimexpr-\dp0-2pt\relax}{\tikz[baseline=0pt,x=1pt,y=1pt]{\draw[line width=.4pt] (0,3)--(0,0)--(\the\wd0,0)--(\the\wd0,3);}}}\copy0\endgroup}<+\infty$.
\margin{117}
